(1) 芯包界面从芯层入射，临界角相对法线满足 \(\sin i_c=n_2/n_1\)。子午光线与光纤轴成角 \(\phi\) 时，在侧壁的入射角 \(i=\pi/2-\phi\)，导波条件 \(i\ge i_c\)，等价于
\[
\sin\phi\le \sqrt{1-(n_2/n_1)^2}.
\]

(2) 光从外界经平面端面进入芯层，Snell 定律 \(n_0\sin\theta=n_1\sin\phi\)。结合最大芯内角，得
\[
n_0\sin\theta_{\max}=\sqrt{n_1^2-n_2^2}.
\]
若右边大于 n0，实际最大接收角受 \(\sin\theta\le1\) 限制，不能形式上把接收角写成大于 90 度。

(3) 定义 \(\mathrm{NA}=\sqrt{n_1^2-n_2^2}\)（通常指外界为空气时的归一化孔径）。导波归一化频率为
\[
V=\frac{2\pi a}{\lambda}\sqrt{n_1^2-n_2^2}.
\]
弱导条件 \(n_1-n_2\ll n_1\) 可使矢量电磁模近似为标量 LP 模。

(4) 当 V≫1 且弱导、阶跃芯层足够均匀时，光纤中可传播模式数近似 \(M\simeq V^2/2\)，这来自横向 k 空间的模态计数；在低 V、强导或光纤弯曲明显时，几何光学估计不准确，应直接求矢量 Maxwell 本征模。
